\scanpage{84}
\uline{\unsure{Modulis|Modulus} $M_{0,4}$ et l'octaèdre}

$I$ ens de card. $4$ d'où

$J=\mathscr P_{2,2}(I)$, $V(J)$ plan vectoriel sur $\mathbb F_2$ avec $V(J)^*\simeq J$,

$V(J)\subset\unsure{\mathfrak S|\mathfrak G}_I$ formé de l'identité et des ``bitranspositions''.

On s'intéresse aux droites projectives (sur $\mathbb C$ disons -- ça marcherait sur tt corps alg clos \inserted{de car. $\ne2$}) munies de $I\hookrightarrow X$. (plus précisément $I_{\mathbb C}\hookrightarrow X$, ou $I\to X(\mathbb C)$ ...).

L'opération de $V(J)$ sur $I$ se prolonge en une opération de $V(J)$ sur $X$
\[
V(J)\hookrightarrow G_X=\operatorname{Aut}(X)
\]

Soit \qquad $S=$ ens des pts de $X$ qui sont fixes par un élément de $J=V(J)^*$
\[
\text{d'où}\qquad \operatorname{card} S=6
\]

% slanted
\margin{\uline{\unsure{Commencer}} \uline{par Ça !}}
En fait, on trouve que \deleted{\unsure{pour tt}} \inserted{les} opérations \deleted{de} \inserted{fidèles} d'un plan vectoriel $V$ sur un $X$, i.e les plongements $V\deleted{\unsure{(J)}}\subset G_X$, sont tous équivalents (modulo autom intérieurs de $G_X$) et sont (pour une base fixée \inserted{$e_0,e_1$}) équivalents à l'opération sur $\mathbb P^1$
\[
\begin{aligned}
u(e_0)&=(z\longmapsto-z)=\begin{pmatrix}-1&0\\0&1\end{pmatrix}&&\text{pts fixes }0,\infty\\
u(e_1)&=(z\longmapsto1/z)=\begin{pmatrix}0&1\\1&0\end{pmatrix}&&\text{pts fixes }\pm1\\
u(e_0+e_1)&=(z\longmapsto-1/z)=\begin{pmatrix}0&-1\\1&0\end{pmatrix}&&\text{pts fixes }\pm i
\end{aligned}
\]

Les \deleted{\unsure{groupes gén}} \inserted{sous groupes} isom. à $D_2$ de $G_X$ correspondent $1$-$1$ aux parties $S$ de cardinal $6$ de $X$, isom. à $(\mathbb P_1,\{0,\infty,\pm1,\pm i\})$.
% slanted
\margin{parties octaédrales d'une droite projective}

Le normalisateur d'un tel \deleted{$D_2$} sous groupe $V$ \inserted{dans $G_X$} est isomorphe au groupe \deleted{des autom} $\mathfrak S_4$ (de façon non canonique), plus précisément, si $V=V(J)$ avec $J=\unsure{J|V|\mathfrak P}(I)$, \inserted{i.e. si $I$ est un $V$-torseur}, on trouve \inserted{alors} iso de suites exactes
\diagram{\begin{tikzcd}[ampersand replacement=\&]
\& \& \deleted{\unsure{\operatorname{Aff}(V)}} \arrow[d,equal] \& \& \\
1\arrow[r] \& V\arrow[r]\arrow[d,equal] \& \mathfrak S_I\arrow[r]\arrow[d] \& \operatorname{Aut}(V)\arrow[r]\arrow[d,equal] \& 1\\
1\arrow[r] \& V\arrow[r] \& \Gamma\arrow[r] \& \operatorname{Aut}(V)\arrow[r,dash] \& 1
\end{tikzcd}}
défini modulo autom intérieurs par un élément de $V$. Quand $I$ est le torseur trivial, se donner un tel \deleted{\unsure{iso}} iso d'ext.
% archive sheet number: 86
